\documentclass[a4paper,12pt]{article}

\usepackage[utf8]{inputenc}
\usepackage[T1]{fontenc}
\usepackage[english]{babel}
\usepackage{amsmath}
\usepackage{amssymb}
\usepackage{graphicx}
\usepackage{hyperref}
\usepackage{geometry}
\usepackage{fontspec}
\usepackage{booktabs}
\usepackage{amsthm}
\usepackage{enumitem}
\usepackage{caption}
\usepackage{thmtools}

% --- PACCHETTI AGGIUNTI PER LA FORMATTAZIONE ---
\usepackage{parskip}  % Aggiunge spazio tra i paragrafi (stile "blocco")
\usepackage{fancyhdr} % Per intestazioni e piè di pagina personalizzati

% --- IMPOSTAZIONI GEOMETRIA E FONT ---
% \geometry{margin=2.5cm} % Questa riga è duplicata e sovrascritta
\geometry{a4paper, margin=1in} % Manteniamo questa, come nell'esempio


% Configura lo stile 'plain' (usato da \maketitle e \tableofcontents)
% per essere coerente con il piè di pagina
\fancypagestyle{plain}{
    \fancyhf{} % Pulisce
    \cfoot{\thepage} % Mette solo il numero di pagina
    \renewcommand{\headrulewidth}{0pt} % Niente linea su queste pagine
}


\begin{document}

\section*{Peer Review C3 - Problem 6.3}

Given that $2$ is a primitive root modulo $101$, find $x \in \mathbb{Z}/101\mathbb{Z}$ such that $ord_{101}(x)=10$.

---

\section*{Solution}

The modulus is $p=101$, which is a prime number.
The primitive root (admitted) is $g=2$.

The order of the multiplicative group $\mathbb{Z}/101\mathbb{Z}^{\times}$ is $\phi(101) = 101-1 = 100$.
Every element $x \in \mathbb{Z}/101\mathbb{Z}^{\times}$ can be written as a power of the primitive root $g$, i.e., $x = g^i = 2^i$ for some integer $i$.

The order of a power $g^i$ is given by the formula:
$$ord_{p}(g^{i})=\frac{p-1}{\gcd(p-1,i)}$$

In our case, we have $p=101$ and we require $ord_{101}(x)=10$. Substituting the known values:
$$10 = \frac{100}{\gcd(100,i)}$$

This implies that $\gcd(100,i)$ must be $10$:
$$\gcd(100,i) = \frac{100}{10} = 10$$

We seek an integer $i$ such that $\gcd(100, i) = 10$.
Let $i = 10k$ for some integer $k$. Substituting this:
$$\gcd(100, 10k) = 10$$
Factoring out $10$:
$$10 \cdot \gcd(10, k) = 10$$
Dividing by $10$:
$$\gcd(10, k) = 1$$

We have to choose an integer $k$ such that $k$ is coprime to $10$. The easiest number we can choose is $k=1$, because $\gcd(10, 1) = 1$.

Thus, we set $k=1$, which gives the exponent $i$:
$$i = 10k = 10 \cdot 1 = 10$$

The element $x$ we are looking for is $x = 2^i$:
$$x = 2^{10} \pmod{101}$$

We calculate the value:
$$2^{10} = 1024$$
Now, we find $1024 \pmod{101}$:
$$1024 \equiv 14 \pmod{101}$$
$$x=14$$

The required value is $\mathbf{x=14}$.

\end{document}